% Generated draft: requires mathematical and source review.
\documentclass[11pt]{article}
\usepackage[margin=25mm]{geometry}
\usepackage{fontspec}
\setmainfont{Latin Modern Roman}
\usepackage{amsmath,amssymb,mathtools,unicode-math}
\setmathfont{Latin Modern Math}
\usepackage{enumitem,needspace,xurl,hyperref}
\setlength{\parindent}{0pt}
\setlength{\parskip}{5pt}
\setlength{\emergencystretch}{3em}
\newcommand{\N}{\mathbb{N}}
\newcommand{\Z}{\mathbb{Z}}
\newcommand{\Q}{\mathbb{Q}}
\newcommand{\R}{\mathbb{R}}
\newcommand{\C}{\mathbb{C}}
\newcommand{\F}{\mathbb{F}}
\newcommand{\A}{\mathbb{A}}
\newcommand{\PP}{\mathbb{P}}
\newcommand{\E}{\mathbb{E}}
\newcommand{\eps}{\varepsilon}
\newcommand{\dd}{\,\mathrm{d}}
\begin{document}
\section*{UnsolvedMath ID: 20001979}
\textbf{Source problem code:} catalogue entry
\subsection{Definitions and mathematical statement}
Let \(M\) be a compact smooth manifold, let \(i:M\hookrightarrow T^*M\) be its zero section, and let \(\sigma\) be the canonical symplectic form. M. Libine's Question 11.1 asks whether, for every exact symplectic form \(\omega=\dd\alpha\) on \(T^*M\) with \(i^*\alpha\) exact, there exist a diffeomorphism \(\Phi:T^*M\to T^*M\) and \(\varepsilon\in\{1,-1\}\) such that
\[
 \Phi(i(M))=i(M),\qquad \Phi^*\omega=\varepsilon\sigma.
\]
Using the inverse diffeomorphism gives the equivalent convention for ``sending'' forms. Preservation of the zero section is taken setwise; no asymptotic or completeness condition is stated in the source.

The supplied aid reports a counterexample to this literal assertion. Take
\[
 M=S^1=\R/(2\pi\Z),\qquad T^*S^1\cong S^1\times\R,
 \qquad \sigma=\dd p\wedge\dd q,
\]
and define
\[
 \omega=e^{-p^2}\dd p\wedge\dd q,
 \qquad
 \alpha=\left(\int_0^p e^{-s^2}\dd s\right)\dd q.
\]
Then \(\omega=\dd\alpha\) is symplectic and \(i^*\alpha=0\). With unsigned symplectic area,
\[
 \int_{T^*S^1}|\omega|=2\pi\sqrt\pi<\infty,
 \qquad \int_{T^*S^1}|\sigma|=\infty.
\]
The aid uses diffeomorphism invariance of the integral of a positive density to exclude equivalence to either sign of \(\sigma\), even without the zero-section-preservation requirement.

It also reports the family \(\omega_c=c\omega\), \(c>0\), with primitives \(c\alpha\). Its members have distinct total areas \(2\pi c\sqrt\pi\), so are pairwise inequivalent under diffeomorphisms pulling one form back to either sign of another. For an equivalence that preserves the zero section, the unordered pair
\[
 \left\{\int_{\{p>0\}}|\omega|,\ \int_{\{p<0\}}|\omega|\right\}
\]
is likewise invariant, because the two components of the complement can be exchanged. For the displayed Gaussian form each entry equals \(\pi\sqrt\pi\). These are the supplied counterexample and invariants, not a classification of exact cotangent symplectic structures.
\subsection{Short English statement}
The source asks whether every exact cotangent symplectic form with an exact zero-section primitive is standard up to sign and a zero-section-preserving diffeomorphism. The supplied Gaussian form on the cotangent cylinder has finite total area, whereas the canonical form has infinite area, so it refutes the literal assertion. Its positive scalings give the reported continuum of inequivalent examples.
\subsection{Sources}
\begin{itemize}
\item Local UnsolvedMath registry, permanent ID 20001979, exact source code catalogue entry; canonical \texttt{aim-geometry-notes.json}, \texttt{notes[316]}, tag \texttt{question}; AIM Workshop Problem Lists, Geometry, Question 11.1 (M. Libine). Created and updated 14 August 2026.
\item AIM \emph{Moment maps and surjectivity in various geometries}; \url{https://aimath.org/WWN/momentmaps/momentmaps.pdf}. The supplied verification reports that the original PDF was checked; that is a source claim, and the PDF was not accessed here.
\item Supplied attempt 1: result status \texttt{counterexample}, problem status at run \texttt{refuted}, rigor label \texttt{proved}, novelty confidence \texttt{low}; classification \texttt{SOLVED-IN-LITERATURE}, registry status \texttt{solved}. Candidate: Gaussian counterexample family and end-area invariant. No separate literature citation for that classification or example is supplied.
\item Ulam AI / UnsolvedMath Contributors; CC BY 4.0. Machine-generated research aid, not peer reviewed; independently verify before citation.
\end{itemize}
\end{document}
